% !TeX root = ../resume.tex
%-------------------------------------------------------------------------------
%	SECTION TITLE
%-------------------------------------------------------------------------------
\cvsection{Проекты}


%-------------------------------------------------------------------------------
%	CONTENT
%-------------------------------------------------------------------------------
\begin{cventries}

%---------------------------------------------------------
\cventry
  {Яндекс лицей} % Organisation
  {\href{https://github.com/mistmake/Recipes_collection.git}{Десктопное приложение "Сборник рецептов"~\faGithubSquare}} % Project
  {Россия} % Location
  {Осень 2022} % Date(s)
  {
    \begin{cvitems}
      \item {Создал сервис, который парсит сайт с рецептами и выводит их в удобном формате в десктопном приложении.}
      \item {Функции и стек: PyQt для десктопного интерфейса; поиск по блюдам и продуктам с просмотром КБЖУ; SQLite-БД с несколькими таблицами; два парсера для продуктов и блюд, наполняющих базу на \textasciitilde{}6900 продуктов и 1521 блюдо.}
    \end{cvitems}
  }
  \cventry
  {} % Organisation
  {\href{https://github.com/mistmake/Reading_tracker}{Веб-приложение "Reading tracker"~\faGithubSquare}}% Project
  {Россия} % Location
  {Зима 2023} % Date(s)
  {
    \begin{cvitems}
      \item {Разработал Django-приложение для чтения и трекинга EPUB-книг.}
      \item {Есть загрузка файлов .epub, и правильное форматирование для вывода на сайт}
      \item {Есть кастомная модель пользователя с логином, регистрацией, login\textbackslash logout, password reset и связка пользователя с книгами.}
    \end{cvitems}
  }
  \cventry
  {HSE University} % Organisation
  {\href{https://github.com/mistmake/HSE-project-Mini-GPT-cpp.git}{HSE MiniGPT in C++~\faGithubSquare}} % Project
  {Россия} % Location
  {Февраль 2026} % Date(s)
  {
    \begin{cvitems}
      \item {Университетский проект, сделанный в команде на C++. Наша цель была реализовать GPT-like модель, которая может предсказать следующие слова и символы по контексту}
      \item {В процессе разработки были созданы: простой биграм для сравнения с финальной моделью, токенизатор, чтобы показать, как анализируется датасет в продакшене, и сама финальная модель.}
      \item {Мы тренировали модель на датасете tinystories (примерно 2млн символов).}
      \item {Параметры модели: 10,788,929 параметров, embedding dimension 384, 6 transformer blocks, 6 attention heads, training time \textasciitilde{}5 hours.}
      \item {Я отвечал за оркестрацию проекта и создание архитектуры модели. Более того, настройка cmakelists.txt была тоже на мне.}
    \end{cvitems}
  }




    
%---------------------------------------------------------

%---------------------------------------------------------
%---------------------------------------------------------

%---------------------------------------------------------
\end{cventries}
